\Needspace{12\baselineskip}
\begin{corollary}\label{cor:smooth-dual-exact-equation}\index{Gauduchon duality!smooth maximizing test}\label{cor:smooth-dual-threshold-rigidity}
Let $0\leq f\in L^1(\mathrm{d}V)$ be nonzero with $b=b_\omega(f)>0$. If $D_{f^{1/n}}$ has a smooth strictly positive maximizer $T\in\mathscr G_\omega$, then $f^{1/n}$ agrees almost everywhere with a smooth nonnegative function, and every threshold subsolution $u$ is smooth and satisfies
\[
 \mathrm{MA}_\omega(u)=bf\,\mathrm{d}V.
\]
In particular a smooth threshold solution exists, and any two threshold subsolutions differ by a constant.
\end{corollary}
\Needspace{8\baselineskip}
\begin{proof}
Let $u$ be any threshold subsolution; one exists by \cref{thm:individual-density-subsolution}. Write $T=\gamma^{n-1}$, where $\gamma$ is the smooth Gauduchon metric given by the positive $(n-1)$st root of $T$, and set
\[
 a=\left(\frac{bf\,\mathrm{d}V}{\gamma^n}\right)^{1/n}.
\]
Put $A=\omega+\ddc u\geq0$. Since $T$ maximizes, the positive measure $R_{u,T}$ of \cref{prop:threshold-trace-saturation} has mass zero, which gives the equality
\[
 A\wedge\gamma^{n-1}=a\gamma^n.
\]
The coefficient measures of the positive current $A$ are controlled by this trace measure, which is absolutely continuous. Thus $A$ has a matrix density that is positive semidefinite almost everywhere, with coefficients in $L^n$ because $a\in L^n$.

\Needspace{5\baselineskip}
For every smooth Hermitian form $\beta$, \cref{lem:unbounded-gauduchon-trace-density} also gives
\[
 A\wedge\beta^{n-1}
 \geq\left(\frac{bf\,\mathrm{d}V}{\beta^n}\right)^{1/n}\beta^n.
\]
The resulting pointwise matrix inequalities hold almost everywhere, simultaneously for all positive definite matrices: on each coordinate chart, first use a countable dense family of constant positive matrices, realized on smaller charts by global smooth Hermitian forms, and then use continuity in the matrix. The variational formula for $(\det P)^{1/n}$ in the proof of \cref{lem:weak-gaud-determinant-dual} therefore yields
\[
 \left(\frac{A^n}{\gamma^n}\right)^{1/n}\geq a
 \quad\text{almost everywhere}.
\]
Here $A^n$ denotes only the determinant of the measurable coefficient matrix of $A$; no nonlinear product of currents is used. The trace equality gives the reverse inequality, by the arithmetic--geometric mean inequality. Equality then forces all eigenvalues of $A$ relative to $\gamma$ to equal $a$, also where $a=0$. Consequently $A=a\gamma$ as currents.

By \cref{lem:smooth-conformal-current}, applied with $\theta=\omega$, both $a$ and $u$ are smooth, and $f^{1/n}=b^{-1/n}a(\gamma^n/\mathrm{d}V)^{1/n}$ has a smooth nonnegative representative. Since $u$ is smooth, we may now take the top power: $(\omega+\ddc u)^n=a^n\gamma^n=bf\,\mathrm{d}V$. Finally, the function $a$ was defined from $f$, $b$ and $\gamma$ alone, so two threshold subsolutions $u_1,u_2$ satisfy $\omega+\ddc u_1=a\gamma=\omega+\ddc u_2$. Their smooth difference is therefore pluriharmonic, hence constant on the connected compact manifold $X$.
\end{proof}

If $f$ is smooth and strictly positive, \cref{thm:gaud-volume-dual-formula} provides a smooth strictly positive maximizing Gauduchon form, obtained from the smooth equation. The corollary therefore applies, and shows in particular that all threshold subsolutions, bounded or not, differ from the smooth solution by constants.
